\section{Related Work}
\textbf{Foundational Scaling Laws of Neural Language Models.} Foundational scaling studies show language modeling loss follows smooth power-laws in model size $N$, data $D$, and compute $C$ \citep{kaplan2020scaling}, with compute-optimal training prescribing near lockstep growth of parameters and tokens under fixed FLOPs \citep{hoffmann2022training}. Later analyses attribute earlier discrepancies to embedding/non-embedding parameter accounting, last-layer costs, optimizer warmup, and scale-sensitive hyperparameters \citep{pearce2024reconciling,porian2024resolving}, while data-centric refinements examine pruning efficiency \citep{sorscher2022beyond}, repetition effects \citep{hernandez2022scaling}, gzip-based complexity predictors \citep{pandey2024gzip}, constrained or synthetic regimes \citep{muennighoff2023scaling,qin2025scaling}, and task transfer (e.g., translation) \citep{isik2024scaling}. Test-time compute amplification supplies an inference analogue to classical training laws \citep{snell2024scaling}. 

\textbf{RL post-training in LLMs.} In RL, power-law trends similarly link capacity, interaction compute, and performance \citep{hilton2023scaling}; scaling RFT across horizon and compute improves mathematical and coding reasoning \citep{guo2025deepseekr1,kimi2025scaling,mai2025agent,zhang2025landscapeagenticreinforcementlearning,zhang2025surveyreinforcementlearninglarge}, while extended schedules \citep{liu2025prorl}, ultra-low-shot or single-example RL \citep{wang2025reinforcement}, and minimal-data efficiency paradigms \citep{li2025limr} probe data–compute tradeoffs. Instability and uneven gains highlight fragile optimization \citep{zeng2025simplerlzoo,yue2025doesrlincentivize}, and multi-domain mixtures reveal both synergy and interference across math, code, and logic \citep{li2025onedomainhelps,cheng2025revisiting}. Recent work further advances RL-based reasoning through attention-guided process supervision~\citep{nie2026attnpo}, self-search reinforcement learning~\citep{fan2025ssrl}, and structured analysis of LLM reasoning behaviors~\citep{yuan2025understandingllmreasoningabstractive}.

\textbf{Mathematical Reasoning with LLMs.} Mathematical reasoning amplifies these dynamics: accuracy generally scales upward while verification behaviors remain inconsistent \citep{touvron2023scaling}; corpus volume and quality jointly shape attainable curves \citep{ye2024skywork}; multi-task math-generalist training diverges from specialist scaling trajectories \citep{yue2023mammoth}; and RL with code execution induces additional behaviors such as emergent tool use concentrated in math problem solving \citep{zeng2025agent}. Collectively, evidence indicates that reasoning performance is governed by interacting axes of model size, data distribution/quality, training (supervised vs.\ RL) paradigm, and allocation of both training and inference compute, while unified laws for mathematical reasoning remain only partially characterized.